\subsection{変数係数の一次型}\label{節：変数係数の一次型}
倍率と付加項が添字に依存する式を,\sectionref*{節：階比型}の積と\sectionref*{節：階差等比複合型}の引き算で解く.

\subsubsection{例題：倍率をそろえる}
まず
\BookMathLayout{\[
    a_1=2,\qquad
    a_{n+1}=\frac{n+2}{n}a_n+(n+1)(n+2)
    \quad(n\ge1)
  \]}{\[
    \begin{aligned}
      a_1&=2,\\
      a_{n+1}&=\frac{n+2}{n}a_n+(n+1)(n+2)\\
      &\quad(n\ge1)
    \end{aligned}
  \]}
を考える.
倍率の分母$n$と分子$n+2$の間には$n+1$がある.
$n(n+1)$と一つ先の$(n+1)(n+2)$を比べれば,共通する$n+1$が消えて元の倍率を作れる.
\[
  \frac{(n+1)(n+2)}{n(n+1)}=\frac{n+2}{n}
\]
から,\textbf{隣り合う比が元の倍率になる量}$h_n=n(n+1)$を選ぶ(\sectionref*{節：補助数列・正規化}).
$n\ge1$で非零なので,$b_n=\dfrac{a_n}{h_n}$と置いて割ると,
\[
  \frac{a_{n+1}}{(n+1)(n+2)}
  =\frac{a_n}{n(n+1)}+1
\]
となる.従って$b_1=1$, $b_{n+1}=b_n+1$より$b_n=n$であり,
\[
  a_n=n^2(n+1)
\]
を得る.

\bookstep{初項を変える}
初項を一般の$a_1=s$に変えると,$b_1=\dfrac{s}{2}$となる.
同じ更新から$b_n=\dfrac{s}{2}+n-1$なので,
\[
  a_n=n(n+1)\left(n+\frac{s-2}{2}\right)
\]
となる.

\subsubsection{一般の式への拡張}
$a_{n+1}=p_na_n+q_n$ $(n\ge1)$に対して,
$h_1=1$, $h_{n+1}=p_nh_n$とする.
すべての$p_n$が非零なら,$h_n$も非零であり,
\[
  h_n=\prod_{k=1}^{n-1}p_k
\]
である.\footnote{空積は$1$とする.}元の式を$h_{n+1}$で割ると,
\[
  b_n=\frac{a_n}{h_n},\qquad
  b_{n+1}-b_n=\frac{q_n}{h_{n+1}}
\]
となる.$b_1=a_1$から\bookterm{difference}{階差}を足すと,
\[
  b_n=a_1+\sum_{k=1}^{n-1}\frac{q_k}{h_{k+1}}
\]
となり,$a_n=h_nb_n$へ戻して,
\[
  a_n=h_n\left(a_1+\sum_{k=1}^{n-1}\frac{q_k}{h_{k+1}}\right)
\]
を得る.\footnote{先ほどは定数倍して$h_n=n(n+1)$を選んだが,
  その場合の$b_1$を正しく合わせれば同じ答えになる.}
\begin{bookpoint}
  \textbf{\bookterm{normalize}{正規化}の役割}\par
  倍率を取り除くと,残った付加項を足す問題になる.
  積・和の表示を得ることと,最後まで簡単な式で計算できることは別である.
\end{bookpoint}
もし$p_j=0$なら,$a_{j+1}=q_j$はそれ以前の項によらず決まる.
零となる更新で区間を分け,非零の倍率が続く範囲で割る方法を使う.

\subsubsection{別解：同じ更新を引く}
同じ例で,正規化から見つけた$v_n=n^2(n+1)$は元の更新を満たす.
今度は$d_n=a_n-v_n$と引けば,
\[
  d_{n+1}=\frac{n+2}{n}d_n
\]
となり,\bookterm{ratio}{階比}型へ戻る.初項$a_1=s$なら$d_1=s-2$であり,
\[
  d_n=\frac{s-2}{2}n(n+1),\qquad a_n=v_n+d_n
\]
から同じ答えを得る.一つの解$v_n$と\bookterm{homogeneous}{同次}の解の定数倍に分かれた(\sectionref*{節：階差等比複合型}).

\begin{bookpoint}
  \textbf{引くか,割るか}\par
  隣り合う比から割る因子を探すか,同じ更新をする量を引く.
  どちらも,付加項と倍率を分ける操作である.
\end{bookpoint}
操作だけの比較は\sectionref*{節：操作を選ぶ練習},
\bookterm{general}{一般項}まで求める練習は\bookproblemref{practice-connections}{1}{接続演習問1}へ進もう.
